% TOP_PROBLEM: {"id": 30006678, "title": "NP versus P/poly", "category_id": 15, "category": {"id": 15, "name": "computer_science", "display_name": "Computer Science", "description": "Computational complexity, algorithms, and theoretical CS.", "slug": "computer-science", "order_index": 15, "created_at": "2026-07-31T15:26:25.671Z"}, "status": "open"}
% FIRST_POSTED: null
% !TeX program = lualatex
% Compile twice: lualatex open_problems_500.tex
% All references and prose are embedded; no BibTeX database or external inputs.
\documentclass[11pt,a4paper]{article}
\usepackage[margin=24mm,headheight=14pt]{geometry}
\usepackage{fontspec}
\setmainfont{Latin Modern Roman}
\setsansfont{Latin Modern Sans}
\setmonofont{Latin Modern Mono}
\usepackage{amsmath,amssymb,mathtools}
\usepackage{unicode-math}
\setmathfont{Latin Modern Math}
\usepackage{microtype}
\usepackage{enumitem,xurl,needspace}
\usepackage[dvipsnames]{xcolor}
\usepackage{hyperref}
\hypersetup{unicode=true,colorlinks=true,linkcolor=MidnightBlue,urlcolor=MidnightBlue,
 pdftitle={Mathematical Research Notebook},pdfauthor={Source-annotated compilation},bookmarksnumbered=true}
\usepackage{bookmark}
\usepackage{tocloft}
\setlength{\cftsecnumwidth}{3em}
\cftsetpnumwidth{2.5em}
\cftsetrmarg{4em}
\usepackage{titlesec}
\titleformat{\section}{\Large\bfseries\raggedright}{\thesection}{0.7em}{}
\usepackage{fancyhdr}
\pagestyle{fancy}\fancyhf{}
\fancyhead[L]{\small\sffamily Open Problems Math | Research notebook}
\fancyhead[R]{\small 27 September 2026}
\fancyfoot[C]{\thepage}
\setlength{\parindent}{0pt}
\setlength{\parskip}{4pt plus 1pt}
\setlength{\emergencystretch}{3em}
\setcounter{tocdepth}{1}
\setcounter{secnumdepth}{1}
\setlist[itemize]{leftmargin=3em,itemsep=3pt,topsep=4pt,parsep=0pt}
\allowdisplaybreaks
\newcommand{\N}{\mathbb{N}}
\newcommand{\Z}{\mathbb{Z}}
\newcommand{\Q}{\mathbb{Q}}
\newcommand{\R}{\mathbb{R}}
\newcommand{\C}{\mathbb{C}}
\newcommand{\F}{\mathbb{F}}
\newcommand{\A}{\mathbb{A}}
\newcommand{\PP}{\mathbb P}
\newcommand{\E}{\mathbb{E}}
\newcommand{\eps}{\varepsilon}
\newcommand{\dd}{\,\mathrm d}
\newcommand{\sref}[2]{\hyperlink{source:#1:#2}{[S#2]}}
\newcommand{\eref}[2]{\hyperlink{extra:#1:#2}{[E#2]}}
% Turn the next switch off for a shorter reading copy. The quotations preserve
% every exactTarget field, including qualifications omitted from a short summary.
\newif\ifshowcatalogue
\showcataloguetrue
\newcommand{\cataloguescope}[1]{\ifshowcatalogue
 \paragraph{Exact scope in the supplied catalogue.}
 \begin{quote}\small #1\end{quote}\fi}
\usepackage{etoolbox}
\AtBeginEnvironment{itemize}{\small}
% Standalone entry; original section content is delimited for verification.
\begin{document}
\setcounter{section}{0}
%% BEGIN ORIGINAL SECTION CONTAINER
% ================================================================
% problemId: 30006678
% Canonical problem ID: 30006678
% exactTarget SHA-256: 756528905f208d5143ef277d3690ac1eb7529347feb3e4209149f9bc8796564e
\clearpage
\section*{\texorpdfstring{NP versus P/poly}{NP versus P/poly}}\label{30006678}
\subsection{Definitions and mathematical statement}
A Boolean circuit is a finite directed acyclic computation with binary AND/OR gates, unary NOT gates, input bits, constants, and one output. Define $\mathsf{P/poly}$ by
\[
 L\in\mathsf{P/poly}\ \Longleftrightarrow\ \exists C,k\ \forall n\ge0\ \exists C_n:
 |C_n|\le C(n+1)^k,\quad C_n(x)=1_L(x)\ (|x|=n).
\]
All nodes are counted. No algorithm for generating $C_n$ is required. With $\mathsf{NP}$ defined by polynomially checkable witnesses, determine whether $\mathsf{NP}\subseteq\mathsf{P/poly}$. Different languages may have different size bounds.
\par\smallskip\textit{Formulation and concept references:} \sref{30006678}{1}.
\cataloguescope{Determine the truth of NP subseteq P/poly. Here NP consists of languages over finite binary strings accepted by nondeterministic Turing machines in polynomial time. P/poly consists of languages L for which there exist constants C>=1 and integer k>=1 and Boolean circuits C\_n for all n>=0, of size at most C(n+1)\textasciicircum{}k, deciding membership in L on every length-n input. Circuits are finite acyclic single-output circuits with fan-in-two AND/OR, fan-in-one NOT, input bits and hardwired constants0/1; count all nodes. Each language has its own bound and circuit family, with no uniform generation requirement.}
\subsection{Short English statement}
Can every efficiently verifiable decision problem be solved by small circuits when a different circuit may be chosen separately for each input length?
%% BEGIN RESEARCH INSERTION
\subsection{Research attempt}
\paragraph{Current frontier.}
\textit{Research check: 27 September 2026.}
The nonuniform convention is essential: the circuits need not have an
efficient generator. The Karp--Lipton consequence of
$\mathsf{NP}\subseteq\mathsf{P/poly}$ is a collapse of the polynomial
hierarchy, not a known contradiction; see Arora--Barak, Theorem 6.13.
\sref{30006678}{1}. The 2026 circuit lecture retains this distinction.
\sref{30006678}{2}. Aaronson, Section 6.2.5, explains the conditional natural-proofs
barrier: efficiently recognizable, large truth-table properties useful
against small circuits conflict with suitably strong pseudorandom
function assumptions. It is not an unconditional prohibition on all
circuit lower-bound methods. \eref{30006678}{1}.

\paragraph{Attack route.}
Counting produces hard truth tables but does not put them in $\mathsf{NP}$.
A more target-specific route is to exploit SAT self-reducibility and
force any candidate circuit to fail on a verifiably satisfiable input.
A third route would derive an impossible consequence from the
Karp--Lipton collapse, but no such impossibility is presently available
in this argument. We investigate whether small test sets can certify
the self-reduction route, and distinguish semantic witness circuits
from general truth-table properties.

\paragraph{Attempt: the counting argument and its missing verifier.}
A circuit with at most $s$ nodes can be described by gate types, at most
two predecessors per node, and an output address. With a generous
encoding this uses
\[
 O\bigl(s\log(n+s+2)\bigr)\quad\hbox{bits}. \tag{12.1}
\]
This bound includes input nodes and constants as required in the original
statement. Thus the number of functions on $n$ bits computed by such
circuits is at most $2^{O(s\log(n+s+2))}$. For any fixed $k$ and
$s=n^k$, this is much smaller than the $2^{2^n}$ Boolean functions when
$n$ is sufficiently large. Most truth tables escape the bound.

Select the lexicographically first table escaping size $n^k$ by exhaustive
enumeration. This defines a function for each $n$, but its value on one
input is not accompanied by a polynomial-length witness whose validity
we have shown checkable in polynomial time. Moreover, a separate language
for each fixed $k$ would not automatically give one $\mathsf{NP}$ language
escaping \emph{every} polynomial size bound. Both the explicit membership
condition and the quantifier over exponents remain. Counting alone
supplies neither. The standard counting/circuit framework is
\sref{30006678}{1}, Section 6.1; the failed selection procedure is examined
here rather than assumed efficient.

\paragraph{Attempt: exact witness-circuit reformulation for SAT.}
Fix a conventional padded encoding of CNF formulas with input length
$N$, at most $N$ variables, and a polynomial-size verification circuit.
If SAT has polynomial-size decision circuits, there are polynomial-size
circuits $W_N$ that, on \emph{every satisfiable} formula of length $N$,
output a satisfying assignment. To construct $W_N$, perform the usual
prefix search: test satisfiability after fixing the next variable to
zero; choose zero if possible and one otherwise. There are at most $N$
queries. Substituted formulas have length polynomial in $N$. Hardwire
the SAT decision circuits for all query lengths up to that bound, and
combine them using multiplexers and padded encodings. The total number
of nodes is still polynomial in $N$, since the exponent of the assumed
SAT family is one fixed constant. Correctness follows by induction on
the prefix, keeping a satisfiable residual formula at every step.

Conversely, given such a family $W_N$, compute an assignment and evaluate
the original formula on it. On a satisfiable input verification accepts;
on an unsatisfiable input no assignment can pass. Hence
\[
 \mathrm{SAT}\in\mathsf{P/poly}
 \quad\Longleftrightarrow\quad
 \text{such uniformly correct, nonuniform witness circuits exist.}
 \tag{12.2}
\]
By polynomial-time reductions to SAT, the left side is also equivalent
to $\mathsf{NP}\subseteq\mathsf{P/poly}$. No generator for $W_N$ is
asserted; the word ``uniformly'' in the right side means correctness on
\emph{all inputs at that length}, not circuit-generation uniformity.
One can avoid that word entirely by reading the displayed condition as
$\forall F$ of length $N$, if $F$ is satisfiable then $F(W_N(F))=1$.
This self-reduction is the ingredient used in the Karp--Lipton proof in
\sref{30006678}{1}.

The hoped-for negative route is now precise:
\[
 \begin{gathered}
 \forall C,k\ \exists N\ \forall W\text{ with }|W|\leq C(N+1)^k\\
 \exists F\text{ satisfiable of length }N:\quad F(W(F))=0.
 \end{gathered} \tag{12.3}
\]
The formula may depend on the candidate circuit. Reversing those last
quantifiers would be a serious mistake: one fixed satisfiable formula
always has a small circuit that ignores its input and outputs a
hardwired satisfying assignment. An exhaustive counterexample search
is not a proof that it terminates for every candidate at one common
length. A proved termination statement with the quantifiers (12.3)
\emph{would} suffice, even if the search itself were slow; the missing
point is existence, not the computational cost of writing a proof.

\paragraph{Attempt: an obstruction to small universal test sets.}
The following elementary example rules out a proposed shortcut. Let the
target function on $\{0,1\}^n$ be identically zero, and let the competitor
class include every point indicator
\[
 d_a(x)=\bigwedge_{i=1}^n
 \begin{cases}x_i,&a_i=1,\\ \neg x_i,&a_i=0.\end{cases}
 \qquad(a\in\{0,1\}^n).
\]
Each competitor has at most $3n+2$ nodes under the original convention,
by explicitly counting inputs, negations, AND gates and an optional
constant. If a test set $T$ omits $a$, then $d_a$ agrees with the zero
function on all of $T$ but differs at $a$. Therefore any test set that
certifies equality to this target against every competitor must be the
entire cube. This target is extremely easy, yet its exact test-set
requirement is exponential.

This does not exclude a SAT-specific test construction using more
structure than arbitrary circuit agreement. It does invalidate the
lemma ``a small target circuit always has polynomially many distinguishing
inputs.'' Random testing fares no better in this example: $q$ uniform
queries detect $d_a$ with probability at most $q/2^n$. Thus a compact
list of checks cannot be inferred solely from the existence of a small
candidate circuit.

\paragraph{Stress test: the collapse is not a contradiction.}
For clarity, the known collapse has a direct witness interpretation.
For a true formula $\forall u\exists v\,\phi(u,v)$, the assumed SAT
circuits can be compiled, as above, into a polynomial-size circuit
$W(u)$ choosing a satisfying $v$ for every $u$. Conversely, existence of
such a $W$ implies the original quantified formula. Quantifying its
polynomial-length circuit description first gives
$\exists W\forall u\,\phi(u,W(u))$, a $\Sigma_2^{\mathsf P}$ expression.
This yields the Karp--Lipton collapse described in \sref{30006678}{1},
Theorem 6.13. The argument cannot finish with ``the hierarchy cannot
collapse,'' since that assertion has not been proved here.

A witness-circuit obstruction need not be a large property of all truth
tables, so it might avoid the hypotheses of a natural-proofs barrier.
That observation is not a lower bound or evidence that a particular
construction works. Conversely, citing ordinary one-way-function
existence alone without checking the quantitative pseudorandomness
hypotheses would overstate that barrier. A negative resolution of this
entry would imply $\mathsf P\ne\mathsf{NP}$, since uniform polynomial
time supplies polynomial-size circuits. This is a dependency warning,
not a proof of the harder nonuniform claim.

\paragraph{Outcome.}
\textbf{Outcome: Exploratory approach with unresolved gap.}
The attempt isolates an exact all-input witness-circuit target and proves
that generic small-test-set compression fails even for a constant
function. It also audits the two missing steps in a counting attack:
placing the chosen language in $\mathsf{NP}$ and handling every polynomial
exponent for one language. No superpolynomial SAT circuit lower bound or
polynomial-size SAT circuit family is obtained. The auxiliary lemmas
are not claimed historically new.
%% END RESEARCH INSERTION
\subsection{Sources}
\begin{itemize}\raggedright
\item[\textnormal{[S1]}]\hypertarget{source:30006678:1}{} Arora and Barak, Computational Complexity: A Modern Approach, web draft January8,2007, Definitions6.1–6.3 and Theorem6.13; Aaronson, P ?= NP, Conjecture25.
\url{https://theory.cs.princeton.edu/complexity/book.pdf}.
{\small\textit{Catalogue formulation source.}}
\item[\textnormal{[S2]}]\hypertarget{source:30006678:2}{} Boolean Circuits and P/poly, Cornell CS 6810, February 26, 2026.
\url{https://www.cs.cornell.edu/courses/cs6810/2026sp/lec10.pdf}.
{\small\textit{Catalogue research/status reference; not a proof endorsement.}}
%% BEGIN RESEARCH REFERENCES
\item[\textnormal{[E1]}]\hypertarget{extra:30006678:1}{} Scott Aaronson, \emph{P =? NP}, in John F. Nash Jr. and Michael Th. Rassias (eds.), \emph{Open Problems in Mathematics}, Springer (2016), 1--122, Section 6.2.5.
\url{https://www.scottaaronson.com/papers/pnp.pdf}.
{\small\textit{Added research source. Natural-proofs barrier and its conditional pseudorandomness hypotheses. Consulted 27 September 2026.}}
%% END RESEARCH REFERENCES
\end{itemize}

%% END ORIGINAL SECTION CONTAINER
\end{document}
